\section{Implementation}
\label{sec:implementation}

The algorithm is implemented in Python. We make use
of the \texttt{hmmlearn} library to store Hidden Markov Models, to emit sequences of
observations, and to apply the forward algorithm for likelihood calculations. For the
ZSS algorithm, we employ the \texttt{zss} package to perform the Zhang--Shasha tree
edit distance, augmented with our custom distance metric to tailor the comparison to
our models.

\subsection{Node HMM Model Augmentations}
In contrast to the formal description of the algorithm, the concrete implementation
requires structural adjustments to the Hidden Markov Models. Specifically, each model
is augmented with a single entry state and a final exit state. The reason for the
initial incoming transition is that, while we represent each node
as an HMM, this is a simplification, a node may have several incoming and several outgoing transitions. To account for this, 
these incoming and outgoing transitions are
collected and normalized into one incoming and one outgoing state, with transition probabilities
summed and normalized. This ensures that every sequence begins from a
unique entry point and terminates in a unique exit point. 

As well, in the concrete models there are hidden states that are silent. The forward
algorithm cannot directly handle the potentially infinite paths that arise from such
$\epsilon$-transitions, so these states must be removed. Eliminating silent states while
preserving the transition and emission probabilities of the non-silent states requires
an $O(n^3)$ preprocessing step, where $n$ is the number of states in the model. This
transformation is performed once for each node. 

Alternatively, we could have chosen 
to force every silent state to emit a self-identifying
character. We believe this modification would shift the comparison between HMMs toward
a more structural rather than emission-based distance measure. Further research could
evaluate the effect of this approach. 


Because the ZSS algorithm must also account for insertions and deletions, we require a
consistent rule for comparing a node to its absence. We handle this by comparing the 
observed sequences emitted from the node, to the probability that it emits the empty string. 

\[
D(H, \emptyset) \approx \frac{1}{2n} \left(
    \sum_{i=1}^n \log \frac{p(h_i \mid H)}{p(\emptyset \mid H)}
    + \sum_{i=1}^n \log \frac{p(\emptyset \mid H)}{p(h_i \mid H)}
\right).
\]

Here $\emptyset$ denotes the empty string, and $h_i$ are the observed strings generated from the model $H$.


Additionally, certain nodes in the concrete model cannot be represented internally as
Hidden Markov Models. Specifically, the \texttt{BIF} and \texttt{END} nodes are treated
as structural markers rather than probabilistic models. Any time these nodes are compared 
to like named nodes they will return a log-odds of 0.0 (perfect match) and a default high value
otherwise. For this test we chose -5.0 odds. 




\subsection{Memoization}
To reduce redundant computation, many values are memoized during parsing and model
construction. All sample strings required for label comparison are generated once and
stored, together with the probability that the generating model produced it, given by the forward algorith.

Furthermore, when comparing two CM's with the ZSS algorithm, the distance
between nodes is memoized after the first computation. Subsequent distance measures
are retrieved directly from this cache, eliminating repeated evaluations and improving
overall efficiency.






% BRYAN WRITE IMPLEMENTATION DETAILS HERE